\heading{数学物理方法（下）-第12次作业}


\begin{question}
设本征值问题
\[
\nabla^2\Phi+\lambda\Phi=0,\qquad \Phi|_\Sigma=0
\]
的归一化本征函数为 \(\Phi_k\)，本征值为 \(\lambda_k\)。证明当 \(\lambda=0\) 不是本征值时，Poisson 方程第一类边值问题
\[
\nabla^2u=-f,\qquad u|_\Sigma=0
\]
的解为
\[
u=\sum_k\frac{A_k}{\lambda_k}\Phi_k,\qquad
f=\sum_kA_k\Phi_k.
\]
\begin{solution}
设
\[
u=\sum_ku_k\Phi_k.
\]
这里使用的是 Dirichlet Laplacian 的本征函数完备性：满足同一齐次边界条件的函数可以按 \(\{\Phi_k\}\) 展开。由于题设中
\[
f=\sum_kA_k\Phi_k,
\]
代入 Poisson 方程并逐项比较系数即可。
由
\[
\nabla^2\Phi_k=-\lambda_k\Phi_k
\]
得
\[
\nabla^2u=-\sum_k\lambda_ku_k\Phi_k.
\]
与 \(-f=-\sum_kA_k\Phi_k\) 比较，得
\[
u_k=\frac{A_k}{\lambda_k}.
\]
若 \(\lambda=0\) 不是本征值，则不存在除零问题，公式成立。
\end{solution}
\end{question}

\begin{question}
利用“同时用两组本征函数展开”的方法，求矩形区域 \(0\le x\le a,\ 0\le y\le b\) 内 Poisson 方程
\[
u_{xx}+u_{yy}=-f(x,y)
\]
在边界条件
\[
u|_{x=0}=u|_{x=a}=0,\qquad
u_y|_{y=0}=u_y|_{y=b}=0
\]
下的解。
\begin{solution}
取
\[
X_m=\sin\frac{m\pi x}{a},\qquad
Y_n=\cos\frac{n\pi y}{b},
\]
其中 \(m=1,2,\ldots\)，\(n=0,1,2,\ldots\)。展开
\[
f(x,y)=\sum_{m=1}^{\infty}\sum_{n=0}^{\infty}F_{mn}X_mY_n.
\]
则
\[
u(x,y)=\sum_{m=1}^{\infty}\sum_{n=0}^{\infty}
\frac{F_{mn}}{(m\pi/a)^2+(n\pi/b)^2}
\sin\frac{m\pi x}{a}\cos\frac{n\pi y}{b}.
\]
其中
\[
F_{mn}=\frac{4}{ab(1+\delta_{n0})}
\int_0^a\int_0^b f(x,y)
\sin\frac{m\pi x}{a}\cos\frac{n\pi y}{b}\,dy\,dx.
\]
\end{solution}
\end{question}

\begin{question}
单位球内定解问题
\[
\nabla^2u=0,\qquad x^2+y^2+z^2<1,\qquad
u|_{r=1}=f(z).
\]
在球坐标中选择合适的积分核，对 \(\theta\) 作积分变换，求 \(u\) 的像函数。能否对 \(r\) 作积分变换？
\begin{solution}
因边界只依赖 \(z=\cos\theta\)，问题轴对称。取 Legendre 变换
\[
f(\cos\theta)=\sum_{l=0}^{\infty}a_lP_l(\cos\theta),
\qquad
a_l=\frac{2l+1}{2}\int_{-1}^{1}f(\mu)P_l(\mu)\,d\mu.
\]
球内有界解为
\[
u(r,\theta)=\sum_{l=0}^{\infty}a_lr^lP_l(\cos\theta).
\]
因此像函数为
\[
\hat u_l(r)=a_lr^l.
\]
\(r\) 方向不是有限区间上的自伴本征问题，且 \(r=0\) 是奇点；通常不对 \(r\) 作同类积分变换，而是用有界性选出 \(r^l\)。
\end{solution}
\end{question}

\begin{question}
求下列常微分方程边值问题相应的 Green 函数：
\[
\frac d{dx}\left[(1-x^2)y'\right]=-\phi(x),\qquad
y(-1),y(1)\text{ 有界}.
\]
\begin{solution}
算子
\[
L=\frac d{dx}\left[(1-x^2)\frac d{dx}\right]
\]
有零模常数函数，因此 Green 函数需在零均值子空间上理解。令
\[
L_xG(x,\xi)=-\delta(x-\xi)+\frac12.
\]
加入 \(1/2\) 是为了使右端在 \([-1,1]\) 上积分为零，从而满足 Neumann 型问题的相容条件。对 \(x\ne\xi\)，方程化为
\[
\frac d{dx}\left[(1-x^2)G_x\right]=\frac12.
\]
在 \(\xi\) 的左右两侧分别积分，并要求端点 \(x=\pm1\) 处 \((1-x^2)G_x\) 不产生不可积奇性，可得
\[
G_x(x,\xi)=
\begin{cases}
\dfrac{x+1}{2(1-x^2)},&x<\xi,\\
\dfrac{x-1}{2(1-x^2)},&x>\xi.
\end{cases}
\]
同时在 \(x=\xi\) 处有跳跃条件
\[
\left[(1-x^2)G_x\right]_{\xi-}^{\xi+}=-1,
\]
这正对应右端的 \(-\delta(x-\xi)\)。
积分并取对称规范，可写为
\[
G(x,\xi)=\frac12\ln\frac{1-x_<}{1+x_<}
+C(\xi),
\]
其中 \(x_<=\min(x,\xi)\)，常数由所取零均值条件确定。解存在的相容条件为
\[
\int_{-1}^{1}\phi(x)\,dx=0.
\]
\end{solution}
\end{question}

\begin{question}
求一维和二维无界空间 Helmholtz 方程散射问题的 Green 函数，时间因子取 \(e^{-i\omega t}\)。
\begin{solution}
取出射波条件。方程
\[
(\nabla^2+k^2)G=-\delta
\]
在一维中的 Green 函数为
\[
G_1(x,x')=\frac{i}{2k}e^{ik|x-x'|}.
\]
二维中
\[
G_2(\bm r,\bm r')=\frac{i}{4}H_0^{(1)}(k|\bm r-\bm r'|).
\]
其中 \(H_0^{(1)}\) 的远场渐近正对应出射圆柱波。
\end{solution}
\end{question}

\begin{question}
求二维 Poisson 方程
\[
\nabla^2G(x,y;x',y')=-4\pi\delta(x-x')\delta(y-y')
\]
的通解。
\begin{solution}
二维基本解满足
\[
\nabla^2\ln|\bm r-\bm r'|=2\pi\delta(\bm r-\bm r').
\]
因此一个特解为
\[
G_0(\bm r,\bm r')=-2\ln|\bm r-\bm r'|.
\]
通解为
\[
G(\bm r,\bm r')=-2\ln|\bm r-\bm r'|+H(\bm r,\bm r'),
\]
其中 \(H\) 对 \(\bm r\) 是任意调和函数，由具体边界条件确定。
\end{solution}
\end{question}

\begin{question}
求无界空间二维波动问题
\[
u_{tt}-a^2\nabla^2u=f(x,y,t),\qquad
u|_{t=0}=\psi(x,y),\qquad u_t|_{t=0}=\xi(x,y)
\]
相应的 Green 函数。
\begin{solution}
迟滞 Green 函数满足
\[
G_{tt}-a^2\nabla^2G=\delta(x-x')\delta(y-y')\delta(t-t'),\qquad G=0\quad(t<t').
\]
二维波动方程的迟滞 Green 函数为
\[
G=\frac{H(\tau-R/a)}{2\pi a^2\sqrt{\tau^2-R^2/a^2}},
\qquad
\tau=t-t',\quad R=|\bm r-\bm r'|.
\]
故解可由 Duhamel 原理写成源项卷积加初值贡献。
\end{solution}
\end{question}

